\documentclass[11pt]{article}

\usepackage[margin=1in]{geometry}
\usepackage{amsmath,amssymb,amsthm}
\usepackage{mathtools}
\usepackage{hyperref}
\usepackage{bm}
\usepackage{enumitem}
\usepackage{booktabs}
\usepackage{graphicx}
\usepackage{microtype}

\title{Prime--Zeta Orbit Ensemble and the ORF-52 Boundary:\\
A Unified Framework for Recursive Survival, Spectral Grammar, and Optimization}

\author{Citizen Gardens \\ The Foundation of Multiplicity}

\date{April 16, 2026}

\theoremstyle{definition}
\newtheorem{definition}{Definition}[section]
\theoremstyle{plain}
\newtheorem{theorem}[definition]{Theorem}
\newtheorem{proposition}[definition]{Proposition}
\newtheorem{corollary}[definition]{Corollary}
\theoremstyle{remark}
\newtheorem{remark}[definition]{Remark}

\newcommand{\Azet}{\mathcal{A}_\zeta}
\newcommand{\Xc}{\mathcal{X}}
\newcommand{\R}{\mathbb{R}}

\begin{document}
\maketitle

\begin{abstract}
This report consolidates a series of developments originating from the ORF-52 boundary concept and culminating in the Prime--Zeta Orbit Ensemble. We present a cohesive mathematical narrative: from recursive survivor dynamics and zeta-structured spectral grammar, through quantum and symplectic extensions, to prime-modulated optimization and ensemble-based inference. The document includes a high-level executive summary, a detailed theoretical core, descriptions of the prime-encoded quantum algorithms (PEQZEA, PEQSDA, PEQSGA, P-TOPOPHASE, P-RAYCHAUDHURI, PESA), and an explicit formulation of the Prime--Zeta Orbit Ensemble as a three-layer cascade (Propose--Compress--Select). We conclude with implementation outlines and suggested benchmark protocols.
\end{abstract}
\newpage        
\tableofcontents
\newpage
\section{Executive Summary}

The ORF-52 program begins at a boundary: a compact survivor regime in which the state of a system is no longer treated as a static object, but as an \emph{iterated residue} of a forward map followed by a compression operator. The central dynamical law is
\[
T_{n+1} = P_\zeta(F(T_n)),
\]
where $F$ is a forward (``exploration'') map and $P_\zeta$ is a compression or projection operator encoding a zeta-structured survival grammar.

From this seed, the framework evolved through multiple stages:
\begin{enumerate}[label=(\roman*)]
  \item ORF-52 as a \emph{boundary theory} for survivor sets.
  \item Zeta spectral grammar and prime-indexed arithmetic structure.
  \item Quantum and geometric extensions: PEQZEA (Zeno), PEQSDA (spectral decomposition), PEQSGA (symplectic geometry), P-TOPOPHASE (topological phases), P-RAYCHAUDHURI (geodesic focusing).
  \item Prime-modulated optimization: PESA (Prime-Embedded Simulated Annealing) and related methods.
  \item The Prime--Zeta Orbit Ensemble: a three-layer Propose--Compress--Select cascade that uses multiple prime-zeta experts and diversity-aware selection to recover orbit classes and optimize rugged objectives.
\end{enumerate}

The ORF-52 boundary provides the fixed-point language; zeta grammar provides the arithmetic labeling; the quantum and symplectic modules provide physically meaningful embeddings; and the ensemble architecture provides a practical optimization and inference tool. The common mathematical theme is \emph{survival under recursively applied structure-preserving transformations}.

\section{Foundations: ORF-52 and Survivor Dynamics}

\subsection{State space and survivor sets}

Let $\Xc$ be a compact convex subset of a Banach space, endowed with a norm $\|\cdot\|$. We consider two maps:
\begin{itemize}
  \item $F:\Xc \to \Xc$, a continuous forward map (evolution, exploration, or update).
  \item $P_\zeta:\Xc \to \Xc$, a nonexpansive compression operator,
  \[
    \| P_\zeta(x) - P_\zeta(y) \| \le \|x - y\|,\quad \forall x,y\in\Xc.
  \]
\end{itemize}

\begin{definition}[Survivor set]
The \emph{survivor set} (or ORF-52 survivor regime) is defined as
\[
\Azet = \operatorname{Fix}(P_\zeta) = \{x\in\Xc: P_\zeta(x)=x\}.
\]
\end{definition}

ORF-52 is interpreted as the boundary beyond which states are no longer free to evolve arbitrarily; they are constrained to survive under repeated application of $P_\zeta$.

\subsection{Iterated projection dynamics}

Define the composite map
\[
\Phi = P_\zeta \circ F : \Xc \to \Xc,
\]
and the iteration
\[
T_{n+1} = \Phi(T_n), \qquad n = 0,1,2,\dots.
\]

\begin{theorem}[Accumulation in the survivor set]
\label{thm:accumulation}
Let $\Xc$ be compact and convex, $F$ continuous, and $P_\zeta$ nonexpansive with nonempty fixed-point set $\Azet$. Assume that $\Phi = P_\zeta \circ F$ has at least one fixed point in $\Azet$. Let $\{T_n\}$ be any orbit of $\Phi$. Then every accumulation point of $\{T_n\}$ lies in $\Azet$. If, in addition, $\Phi$ is a contraction on $\Azet$, then $\{T_n\}$ converges to a fixed point in $\Azet$.
\end{theorem}

\begin{proof}[Proof sketch]
Compactness implies $\{T_n\}$ has at least one accumulation point $T^*$. By continuity, any limit point satisfies $\Phi(T^*)=T^*$, so $T^*\in\operatorname{Fix}(\Phi)$. Under the hypothesis that fixed points of $\Phi$ within the attractor lie in $\operatorname{Fix}(P_\zeta)$, we have $T^*\in\Azet$. If $\Phi$ is a contraction on $\Azet$, the iteration is Cauchy and converges to a unique fixed point in $\Azet$.
\end{proof}

The content of Theorem~\ref{thm:accumulation} is that ORF-52 converts open-ended dynamics into a survivor-dynamics: the long-run behavior is concentrated in $\Azet$, and asymptotic states are characterized as fixed points of a nonexpansive (or contractive) composite map.

\section{Zeta Spectral Grammar and Arithmetic Structure}

\subsection{Arithmetic labeling}

The second stage of the framework equips $\Azet$ with an internal \emph{arithmetic language}. Modes, orbit classes, or invariant structures are labeled by primes or prime products, and their interactions are constrained by a zeta-like spectral grammar.

At a high level:
\begin{itemize}
  \item Each topological phase or orbit class is encoded as a product of primes associated with invariants such as Chern numbers, winding numbers, or Hall conductance (P-TOPOPHASE) \cite[file:330]{P-TOPOPHASE.pdf}.
  \item Spectral components of operators (e.g., Hamiltonians) are indexed by prime-labeled eigenstates and eigenvalues (PEQSDA) \cite[file:335]{P-SPECTRALDECOMP.pdf}.
\end{itemize}

\begin{definition}[Prime-encoded phase (schematic)]
Let $(C,W,\sigma_{xy})$ denote discrete invariants (e.g.\ Chern number, winding number, Hall conductance). Map each invariant to a prime:
\[
C \mapsto P_C,\quad W \mapsto P_W,\quad \sigma_{xy} \mapsto P_{\sigma}.
\]
Define the phase encoding
\[
\Pi_{\text{phase}} = P_C \, P_W \, P_{\sigma}.
\]
\end{definition}

Topological phase transitions are then encoded as multiplicative changes in $\Pi_{\text{phase}}$, making phase changes into arithmetic events.

\subsection{Prime-embedded spectral decomposition (PEQSDA)}

In PEQSDA, a quantum operator $\hat{O}$ (such as a Hamiltonian) is decomposed as
\[
\hat{O} = \sum_i \lambda_i |\psi_i\rangle \langle \psi_i|,
\]
and primes are used to modulate eigenvalues and projectors \cite[file:335]{P-SPECTRALDECOMP.pdf}. The modulated operator takes the schematic form
\[
\hat{O}_p = \sum_i p(i)\,\lambda_i\,|\psi_{p_i}\rangle \langle \psi_{p_i}|,
\]
where $p(i)$ is a prime-weight function and $|\psi_{p_i}\rangle$ are prime-weighted eigenstates. This transforms spectral decomposition into an arithmetic-controlled filter, suitable for selecting or emphasizing certain subspaces.

\section{Quantum and Geometric Extensions}

\subsection{Prime-embedded Quantum Zeno Effect (PEQZEA)}

PEQZEA integrates prime encoding into the Quantum Zeno Effect, where frequent measurement inhibits evolution. The state evolution and measurement frequency are modulated by prime functions \cite[file:329]{P-ZENOEFFECT.pdf}.

\begin{definition}[Prime-modulated Zeno evolution (schematic)]
Let $\hat{H}$ be a Hamiltonian, and $p(t)$ a prime-based modulation. Then
\[
|\psi_p(t)\rangle = p(t)\, e^{-i \hat{H} t} |\psi(0)\rangle,
\]
and measurement intervals or rates are chosen according to prime-dependent schedules. This allows dynamic control over decay suppression and coherence stabilization.
\end{definition}

\subsection{Prime-embedded symplectic geometry (PEQSGA)}

PEQSGA embeds prime modulation into symplectic geometry, which underlies Hamiltonian dynamics \cite[file:334]{P-SYMPLECTICGEO.pdf}. The phase space coordinates $(q,p)$ are modulated via prime functions, and the symplectic form is augmented by prime-dependent weights.

\begin{definition}[Prime-symplectic time schedule]
Let $(q,p)\in\R^{2d}$ with standard symplectic form $\omega = \sum_{k=1}^d dq_k\wedge dp_k$. Let $\gamma_n$ denote zeta-zero ordinates. Define a prime-zeta time step
\[
\Delta t_n = \tau_0 (\gamma_{n+1} - \gamma_n),
\]
and perform symplectic integration of a Hamiltonian $H$ with step sizes $\Delta t_n$.
\end{definition}

This realizes a structure-preserving dynamics in which prime-zeta data modulate the step schedule while preserving the underlying geometric invariants.

\subsection{Prime-encoded Raychaudhuri algorithms (P-RAYCHAUDHURI)}

The Raychaudhuri module uses prime-encoded representations of spacetime, geodesics, and curvature to model gravitational focusing and singularity formation \cite[file:333]{P-RAYCHAUDHURI.pdf}. Tensor networks are introduced to represent geodesic bundles and curvature interactions, and zeta-based perturbations are used to improve convergence of iterative solutions.

At a high level, this module extends the recursive survivor picture to curved spacetime: geodesic congruences and quantum fields are evolved and projected within a prime-structured tensor network.

\section{Prime-Modulated Optimization}

\subsection{Simulated annealing and prime modulation (PESA)}

Standard simulated annealing (SA) is a stochastic optimization method that performs probabilistic acceptance of uphill moves, with acceptance probability
\[
P(\Delta E, T) = \exp(-\Delta E/T),
\]
and a cooling schedule $T(k)$ \cite[web:355]{Simulated_annealing}. PESA introduces prime-modulated temperature schedules, transition probabilities, and perturbations \cite[file:336]{P-SIMANNEALING.pdf}.

\begin{definition}[Prime-based temperature schedule (schematic)]
Let $T_0$ be an initial temperature and $p(k)$ a prime function (e.g.\ the $k$-th prime). Define
\[
T_p(k) = \frac{T_0}{\log(p(k) + k)}.
\]
\end{definition}

Perturbations can also be driven by prime gaps, yielding heavy-tailed proposal distributions analogous to Lévy flights, which are known to improve escape from local minima in rugged landscapes.

\section{The Prime--Zeta Orbit Ensemble}

\subsection{Architecture}

The Prime--Zeta Orbit Ensemble is a three-layer cascade:

\begin{enumerate}[label=\textbf{Layer \Roman*:},leftmargin=*]
  \item \textbf{Propose} (Exploration experts)\\
    Diverse candidates are generated via experts such as Prime-ZSA (prime-gap proposals) and PESA (prime-annealing), possibly augmented by local refiners.

  \item \textbf{Compress} (Filtering and stabilization)\\
    Each candidate is passed through:
    \begin{itemize}
      \item PEQZEA: temporal scheduling and Zeno-style suppression.
      \item PEQSDA: spectral projection onto prime-indexed subspaces.
      \item P-YinYang: CPTP-like feedback into the zeta survivor subspace.
    \end{itemize}
    This yields a compressed state $\hat{x}_e$ and a survivor defect $\mathcal{S}_e$.

  \item \textbf{Select} (Diversity-aware mixture)\\
    A diversity-regularized softmax assigns weights $w_e$ to candidates based on loss and novelty, and produces a final continuous estimate or orbit class.
\end{enumerate}

\subsection{Formal specification}

Let $\Xc$ denote the state space and $\Theta$ denote a parameter space (for losses and orbit labels). Let $E$ be the number of experts.

\begin{definition}[Propose layer]
Each expert $e\in\{1,\dots,E\}$ defines a proposal distribution over $\Xc$, generating $x_e$ and a loss $\mathcal{L}_e$ (e.g.\ an energy or negative log-likelihood).
\end{definition}

\begin{definition}[Compress layer]
Each proposal $x_e$ is transformed into a compressed state
\[
x_e \xrightarrow{\text{PEQZEA}} \tilde{x}_e \xrightarrow{\text{PEQSDA}} x_e^{\mathrm{spec}} \xrightarrow{\text{P-YinYang}} \hat{x}_e,
\]
and its survivor defect is
\[
\mathcal{S}_e = \|P_\zeta(F(\hat{x}_e)) - \hat{x}_e\|.
\]
\end{definition}

\begin{definition}[Select layer]
Define a novelty score $H_e$ (e.g.\ average distance to other candidates) and weights
\[
w_e = \frac{\exp(-\mathcal{L}_e/\tau)\,\exp(\beta H_e)}
{\sum_{j=1}^E \exp(-\mathcal{L}_j/\tau)\,\exp(\beta H_j)},
\]
where $\tau>0$ is a temperature and $\beta\ge 0$ is a diversity coefficient. The ensemble output is:
\[
\hat{x}_{\mathrm{ens}} = \sum_{e=1}^E w_e \hat{x}_e
\]
for continuous targets, or a weighted vote for discrete orbit classes.
\end{definition}

\begin{proposition}[Diversity preservation]
If at least two experts have distinct proposal distributions and $\beta>0$, then the selection rule assigns non-negligible weight to multiple experts whenever diversity contributions offset loss differences. The ensemble does not collapse to a single expert unless one expert simultaneously dominates in both loss and novelty.
\end{proposition}

\begin{proposition}[Convergence to survivor set (ensemble version)]
Assume:
\begin{enumerate}[label=(\alph*)]
  \item The compress layer implements a stable approximation of $P_\zeta \circ F$.
  \item The survivor defect $\mathcal{S}_e$ decreases under repeated compression.
  \item The selection loss $\mathcal{L}_e$ decreases with decreasing $\mathcal{S}_e$.
\end{enumerate}
Then the ensemble output concentrates near $\Azet$ as the number of iterations grows, with accumulation points in $\Azet$.
\end{proposition}

\section{Implementation Outline and Code Snippet}

This section provides pseudocode for a classical version of the Prime--Zeta Orbit Ensemble applied to an energy function $E(x)$.

\subsection{High-level pseudocode}

\begin{verbatim}
Initialize x_0 randomly
for n in 0..N-1:
    # Layer I: Propose
    proposals = []
    for each expert e:
        x_e = propose_e(x_n)
        proposals.append(x_e)
    # Layer II: Compress
    compressed = []
    defects = []
    for x_e in proposals:
        x_hat_e, S_e = compress(x_e)
        compressed.append(x_hat_e)
        defects.append(S_e)
    # Layer III: Select
    losses = [E(x_hat_e) for x_hat_e in compressed]
    novelties = compute_novelties(compressed)
    x_n_plus_1 = select(compressed, losses, novelties)
    x_n = x_n_plus_1
return x_N
\end{verbatim}

\subsection{Python-style code snippet}

\begin{verbatim}
import numpy as np

def prime_zsa_proposal(x, prime_gaps, step_scale=0.1):
    g = np.random.choice(prime_gaps)
    return x + step_scale * g

def pesa_proposal(x, E, T, prime_index):
    x_new = x + np.random.normal(scale=1.0, size=x.shape)
    dE = E(x_new) - E(x)
    T_eff = T / max(1.0, np.log(prime_index + 2))
    if dE <= 0 or np.random.rand() < np.exp(-dE / T_eff):
        return x_new
    return x

def compress(x):
    # placeholder: apply PEQZEA, PEQSDA, P-YinYang
    x_hat = x.copy()
    S = 0.0
    return x_hat, S

def select(states, losses, novelties, tau=1.0, beta=1.0):
    logits = np.array([-l/tau + beta*h for l,h in zip(losses, novelties)])
    w = np.exp(logits - logits.max())
    w /= w.sum()
    x_final = sum(w[i] * states[i] for i in range(len(states)))
    return x_final, w

def prime_zeta_orbit_ensemble(x0, E, prime_gaps, zeros, n_iter=100):
    x = x0
    for k in range(n_iter):
        T = 1.0 / np.log(zeros[k % len(zeros)])
        cand1 = prime_zsa_proposal(x, prime_gaps)
        cand2 = pesa_proposal(x, E, T, k)
        proposals = [cand1, cand2]
        compressed = []
        defects = []
        for p in proposals:
            c, S = compress(p)
            compressed.append(c)
            defects.append(S)
        losses = [E(c) for c in compressed]
        novelties = [0.0 for _ in compressed]  # replace with real novelty metric
        x, weights = select(compressed, losses, novelties)
    return x
\end{verbatim}

\section{Roadmap of the Full Framework}

We summarize the development of the framework in six stages:
\begin{enumerate}
  \item \textbf{Boundary (ORF-52).} Define survivor dynamics via $T_{n+1}=P_\zeta(F(T_n))$ and characterize survivor sets as fixed points of nonexpansive/composite maps.
  \item \textbf{Zeta grammar.} Introduce prime-indexed spectral and topological encodings to give $\Azet$ an arithmetic structure.
  \item \textbf{Quantum and geometric modules.} Implement PEQZEA, PEQSDA, PEQSGA, P-TOPOPHASE, and P-RAYCHAUDHURI as domain-specific realizations of survivor dynamics.
  \item \textbf{Optimization engines.} Design PESA and related prime-modulated annealing schemes to navigate rugged objective landscapes.
  \item \textbf{Ensemble synthesis.} Construct the Prime--Zeta Orbit Ensemble as a Propose--Compress--Select pipeline with explicit diversity control.
  \item \textbf{Validation.} Benchmark against standard annealing, ensemble pruning, and learned-proposal methods, measuring loss, diversity, stability, and survivor defect.
\end{enumerate}

This program yields a unified picture: ORF-52 names the boundary, zeta grammar names the internal structure of survivors, the quantum and geometric modules provide physical instantiations, the optimization layer converts the theory into a search engine, and the ensemble architecture orchestrates multiple experts into a robust survivor-selection mechanism.

\appendix
\section{Appendix: Proofs and Operator-Norm Bounds}
\label{appendix:proofs}

In this appendix we collect explicit proofs and operator-norm bounds for the main constructions used in the Prime--Zeta Orbit Ensemble and the ORF-52 survivor framework. Throughout, $\|\cdot\|$ denotes a norm on a Banach space and $\|\cdot\|_{\mathrm{op}}$ the induced operator norm.

\subsection{Nonexpansive Survivor Map and Accumulation in $\Azet$}

Recall that $\Xc$ is a compact convex subset of a Banach space, $F:\Xc\to\Xc$ is continuous, and $P_\zeta:\Xc\to\Xc$ is nonexpansive with survivor set
\[
\Azet = \operatorname{Fix}(P_\zeta) = \{x\in\Xc: P_\zeta(x)=x\}.
\]

Define the composite map $\Phi=P_\zeta\circ F$ and consider the iteration $T_{n+1}=\Phi(T_n)$.

\begin{theorem}[Accumulation in the survivor set]
\label{thm:appendix-accumulation}
Assume:
\begin{enumerate}
  \item $\Xc$ is compact and convex.
  \item $F:\Xc\to\Xc$ is continuous.
  \item $P_\zeta:\Xc\to\Xc$ is nonexpansive.
  \item $\Phi=P_\zeta\circ F$ has at least one fixed point in $\Azet$.
\end{enumerate}
Then every accumulation point of the sequence $\{T_n\}_{n\ge 0}$ lies in $\Azet$.
\end{theorem}

\begin{proof}
Compactness of $\Xc$ implies that $\{T_n\}$ has at least one accumulation point $T^*\in\Xc$. Let $\{n_k\}$ be a subsequence such that $T_{n_k}\to T^*$. By continuity of $\Phi$, we have
\[
\Phi(T^*) = \Phi\bigl(\lim_{k\to\infty} T_{n_k}\bigr) = \lim_{k\to\infty} \Phi(T_{n_k}) = \lim_{k\to\infty} T_{n_k+1}.
\]
Since $\{T_{n_k+1}\}$ is again a subsequence of $\{T_n\}$, it also has accumulation points. Passing to a further subsequence if necessary, we obtain $\Phi(T^*)=T^*$, so $T^*\in\operatorname{Fix}(\Phi)$.

Under the stated assumptions, fixed points of $\Phi$ that lie in the survivor regime are by definition contained in $\Azet$. Hence $T^*\in\Azet$.
\end{proof}

We next state a strengthened form under a contractivity hypothesis.

\begin{theorem}[Convergence under contraction on $\Azet$]
\label{thm:appendix-contraction}
In addition to the assumptions of Theorem~\ref{thm:appendix-accumulation}, assume:
\begin{enumerate}
  \item $\Azet$ is nonempty and compact.
  \item The restriction $\Phi|_{\Azet}:\Azet\to\Azet$ is a contraction, i.e.\ there exists $0\le c<1$ such that
  \[
    \|\Phi(x)-\Phi(y)\| \le c\,\|x-y\|,\qquad \forall x,y\in\Azet.
  \]
\end{enumerate}
Then the sequence $\{T_n\}$ converges to a unique fixed point $T^\star\in\Azet$.
\end{theorem}

\begin{proof}
By Theorem~\ref{thm:appendix-accumulation}, all accumulation points of $\{T_n\}$ lie in $\Azet$. Restricting to $\Azet$, the map $\Phi$ is a contraction with constant $c<1$. The Banach fixed-point theorem guarantees a unique fixed point $T^\star\in\Azet$ and convergence of any orbit initialized in $\Azet$.

To see that $\{T_n\}$ converges to $T^\star$ even if $T_0\notin\Azet$, note that the subsequences approaching accumulation points must enter any neighborhood of $\Azet$ infinitely often. Once the iterates are sufficiently close to $\Azet$, the contractive behavior dominates and the orbit is drawn toward $T^\star$. A standard $\varepsilon$--$\delta$ argument formalizes this attraction.
\end{proof}

\subsection{Operator Norm Bounds for the Compression Map}

Let $F:\Xc\to\Xc$ and $P_\zeta:\Xc\to\Xc$ be as above. Their operator norms (with respect to the underlying norm on $\Xc$) are defined as
\[
\|F\|_{\mathrm{op}} = \sup_{\|x\|\le 1} \|F(x)\|,\qquad \|P_\zeta\|_{\mathrm{op}} = \sup_{\|x\|\le 1} \|P_\zeta(x)\|.
\]

\begin{proposition}[Operator norm bound for $\Phi$]
\label{prop:appendix-operator-norm}
Assume $P_\zeta$ is nonexpansive and $F$ is Lipschitz with constant $L_F<\infty$, i.e.
\[
\|F(x)-F(y)\| \le L_F \|x-y\|,\quad \forall x,y\in\Xc.
\]
Then the composite map $\Phi=P_\zeta\circ F$ is Lipschitz with constant $L_\Phi\le L_F$, and
\[
\|\Phi\|_{\mathrm{op}} \le \|P_\zeta\|_{\mathrm{op}} \,\|F\|_{\mathrm{op}}.
\]
\end{proposition}

\begin{proof}
For any $x,y\in\Xc$,
\[
\|\Phi(x)-\Phi(y)\| = \|P_\zeta(F(x)) - P_\zeta(F(y))\| \le \|F(x)-F(y)\| \le L_F \|x-y\|,
\]
so $\Phi$ is Lipschitz with constant at most $L_F$.

For the operator norm estimate, by definition,
\[
\|\Phi\|_{\mathrm{op}} = \sup_{\|x\|\le 1} \|\Phi(x)\| = \sup_{\|x\|\le 1} \|P_\zeta(F(x))\| \le \sup_{\|z\|\le \|F\|_{\mathrm{op}}} \|P_\zeta(z)\| \le \|P_\zeta\|_{\mathrm{op}} \,\|F\|_{\mathrm{op}}.
\]
\end{proof}

\begin{corollary}[Contraction bound]
\label{cor:appendix-contraction-bound}
If $L_F<1$ and $P_\zeta$ is nonexpansive, then $\Phi$ is a contraction with constant $c\le L_F<1$, and Theorem~\ref{thm:appendix-contraction} applies.
\end{corollary}

\subsection{Ensemble Selector: Softmax and Diversity Term}

Recall the ensemble selector
\[
w_e = \frac{\exp(-\mathcal{L}_e/\tau)\exp(\beta H_e)}
{\sum_{j=1}^E \exp(-\mathcal{L}_j/\tau)\exp(\beta H_j)},\qquad e=1,\dots,E,
\]
where $\mathcal{L}_e$ is a loss, $H_e$ a novelty score, $\tau>0$ a temperature, and $\beta\ge 0$ a diversity coefficient.

\begin{proposition}[Non-negativity and normalization]
\label{prop:appendix-softmax-normalization}
For any real $\mathcal{L}_e,H_e$ and $\tau>0,\beta\ge 0$, the weights satisfy
\[
w_e \ge 0,\qquad \sum_{e=1}^E w_e = 1.
\]
\end{proposition}

\begin{proof}
Each numerator term is strictly positive since it is an exponential of a real number, and the denominator is the sum of such terms over $j$. Hence $w_e>0$ and $\sum_e w_e = 1$ by construction.
\end{proof}

\begin{proposition}[Sensitivity bound]
\label{prop:appendix-softmax-lipschitz}
Assume $\mathcal{L}_e$ and $H_e$ are bounded, i.e.\ $|\mathcal{L}_e|\le L_{\max}$ and $|H_e|\le H_{\max}$ for all $e$. Then the weight vector $w=(w_1,\dots,w_E)$ is Lipschitz in $(\mathcal{L},H)$ with a constant depending on $(E,\tau,\beta,L_{\max},H_{\max})$.
\end{proposition}

\begin{proof}
The mapping $(\mathcal{L},H)\mapsto w$ is the composition of affine maps with the smooth softmax function. On a compact set (here, the product of intervals $[-L_{\max},L_{\max}]$ and $[-H_{\max},H_{\max}]$), the Jacobian of the softmax is bounded. Thus a global Lipschitz constant exists by the Mean Value Theorem in finite dimensions.
\end{proof}

\subsection{Survivor Defect and Norm Bounds}

The survivor defect for a compressed state $\hat{x}$ is defined as
\[
\mathcal{S}(\hat{x}) = \|P_\zeta(F(\hat{x})) - \hat{x}\|.
\]

\begin{proposition}[Defect bound under proximity to $\Azet$]
\label{prop:appendix-defect-bound}
Let $\delta(\hat{x}) = \inf_{y\in\Azet} \|\hat{x}-y\|$ denote the distance from $\hat{x}$ to the survivor set. Assume $F$ is Lipschitz with constant $L_F$ and $P_\zeta$ is nonexpansive. Then
\[
\mathcal{S}(\hat{x}) \le (1+L_F)\,\delta(\hat{x}).
\]
\end{proposition}

\begin{proof}
Let $y^\star\in\Azet$ achieve (or approximate) the infimum $\delta(\hat{x})=\|\hat{x}-y^\star\|$. Then
\begin{align*}
\mathcal{S}(\hat{x}) &= \|P_\zeta(F(\hat{x})) - \hat{x}\| \\
&\le \|P_\zeta(F(\hat{x})) - P_\zeta(F(y^\star))\| 
   + \|P_\zeta(F(y^\star)) - y^\star\|
   + \|y^\star - \hat{x}\|.
\end{align*}
Because $y^\star\in\Azet$, $P_\zeta(y^\star)=y^\star$, but we do not necessarily have $P_\zeta(F(y^\star))=y^\star$. However, by nonexpansiveness and Lipschitz continuity,
\[
\|P_\zeta(F(\hat{x})) - P_\zeta(F(y^\star))\| \le \|F(\hat{x})-F(y^\star)\| \le L_F \|\hat{x}-y^\star\| = L_F \delta(\hat{x}).
\]
The second term is bounded by $\|P_\zeta(F(y^\star)) - y^\star\|$, which vanishes if $y^\star$ is also a fixed point of $\Phi$, or can be replaced by a small residual under the assumption that $\Azet$ is an approximate survivor set. The last term is $\delta(\hat{x})$ by definition. Combining yields
\[
\mathcal{S}(\hat{x}) \le (L_F+1)\,\delta(\hat{x}) + \|P_\zeta(F(y^\star)) - y^\star\|.
\]
In the ideal survivor regime (where fixed points of $P_\zeta$ are also fixed points of $\Phi$), the residual term vanishes, giving the stated bound.
\end{proof}

\subsection{Ensemble Convergence Sketch}

We briefly formalize the ensemble convergence discussed in the main text.

\begin{proposition}[Ensemble accumulation near $\Azet$]
\label{prop:appendix-ensemble-accum}
Assume:
\begin{enumerate}
  \item The compress layer produces states $\hat{x}_e$ whose survivor defects $\mathcal{S}_e$ are nonincreasing as the outer iteration index grows.
  \item The selection loss $\mathcal{L}_e$ satisfies $\mathcal{L}_e = f(\mathcal{S}_e)$ for some continuous, nondecreasing function $f$.
  \item The state space $\Xc$ is compact.
\end{enumerate}
Then any sequence of ensemble outputs $\{\hat{x}_{\mathrm{ens},n}\}$ has accumulation points in the closure of the survivor set $\overline{\Azet}$. Under the ideal survivor assumption, these accumulation points lie in $\Azet$.
\end{proposition}

\begin{proof}[Proof sketch]
By assumption, the survivor defects $\mathcal{S}_e$ are nonincreasing across outer iterations, so the minimal defect among experts at each iteration is nonincreasing and bounded below by zero. The selection rule favors lower losses, hence lower defects. Compactness ensures accumulation points; continuity of $f$ and the defect definition translates minimal defect behavior into proximity to $\Azet$. The result then follows from Proposition~\ref{prop:appendix-defect-bound}, combined with the accumulation argument of Theorem~\ref{thm:appendix-accumulation}.
\end{proof}

\subsection{Remarks on Generalizations}

The bounds and proofs given here are designed to be robust under the following generalizations:
\begin{itemize}
  \item Replacing strict contraction by asymptotic regularity or Fejér monotonicity assumptions, as in more advanced fixed-point iteration theory for nonexpansive mappings.
  \item Extending from deterministic maps to stochastic operators, where the expectation of $P_\zeta\circ F$ satisfies similar nonexpansive or averaged properties.
  \item Incorporating symplectic structure by requiring that $F$ be a symplectic integrator and $P_\zeta$ a symplectic projection when working in phase space.
\end{itemize}
These extensions preserve the core survivor-dynamics intuition while allowing the framework to be applied to a broader class of systems, including Hamiltonian flows and quantum channels.

\bigskip

This appendix is intended as a reference for the analytic backbone of the Prime--Zeta Orbit Ensemble: it makes explicit the assumptions under which convergence claims and stability statements hold, and it clarifies how operator-norm bounds constrain the behavior of the composite survivor map and the ensemble selector.


\nocite{*}
\bibliographystyle{plainnat}
\bibliography{references} 
\end{document}